%% toc: P15 158 158.4,158.5
% Hahn partitions up to signed-null sets, and Stieltjes representations of a finite signed measure (Durrett, Appendix A.4, p. 468).
\chapter{Hahn Partitions and Stieltjes Representations}
\label{sec:p15-w02b-p15-005-hahn-stieltjes}

\Cref{sec:p15-radon-nikodym} proves two decompositions whose
uniqueness is deliberately qualified. A Hahn partition
(\cref{thm:p15-hahn-decomposition}) is fixed only up to signed-null
sets, and a finite signed Borel measure on the line has many expressions
as a difference of two cumulative functions
(\cref{thm:p15-signed-cdf-variation}), only one of which comes from its
Jordan components. This chapter makes both qualifications exact. The
source is Durrett's author version of \emph{Probability: Theory and
Examples}, dated 11 January 2019. On printed page~468 it states that the
Hahn decomposition is not unique, describes every Hahn positive set for
an integral-defined signed measure, and records that two Hahn partitions
differ by signed-null sets; the remark beginning on that page writes the
cumulative function of a finite signed measure as a difference of two
bounded increasing functions \cite[pp.~468--469]{DurrettPTEAuthor2019}.

\emph{Inputs.} A signed measure \(\alpha\) in the one-sided convention
of \cref{def:p15-signed-measure}, with one Hahn partition \((A,B)\);
for the density case, a measure \(\mu\) and an extended real measurable
\(f\) with \(\int f^-\,d\mu<\infty\); for the line, a finite signed
Borel measure and a pair of cumulative functions representing it.
\emph{Conclusions.} The complete list of Hahn partitions and the exact
condition for there to be only one; the source's description of the
Hahn positive sets of an integral-defined signed measure; and the exact
relation between an arbitrary pair of Stieltjes measures representing a
finite signed measure and its Jordan components, with witnesses showing
that both kinds of non-uniqueness occur.

Two parts of this chapter go beyond the source text. The source
illustrates non-uniqueness by its density example and proves that
two partitions disagree only on signed-null sets; the converse, that every
signed-null modification of a Hahn partition is again a Hahn partition,
is proved here. The source states the difference representation of the
cumulative function and its converse; the common-excess description of
all such representations is an extension made here.

\section{Hahn partitions up to signed-null sets}
\label{sec:p15-w02b-p15-005-hahn}

\begin{hbdefinition}{def:p15-w02b-p15-005-hahn-partition}{Hahn partitions and almost-everywhere inclusion}
Let \(\alpha\) be a signed measure on \((\Omega,\mathcal F)\).
A \emph{Hahn partition} for \(\alpha\) is a pair \((A,B)\) of
measurable sets with \(A\cap B=\varnothing\), \(A\cup B=\Omega\),
\(A\) positive and \(B\) negative in the sense of
\cref{def:p15-signed-positive-sets}. Since \(B=A^c\), the pair is
determined by \(A\), which is called a \emph{Hahn positive set}.
\Cref{thm:p15-hahn-decomposition} proves that at least one exists.

For measurable sets \(D,E\), the symmetric difference is
\(D\mathbin{\triangle}E\defeq(D\setminus E)\cup(E\setminus D)\).
For a measure \(\mu\), the statement that \(D\subseteq E\)
\emph{\(\mu\)-almost everywhere} means
\[
                       \mu(D\cap E^c)=0 .
\]
This is the source's convention; it writes the inclusion sign
\(\subset\) for an inclusion that permits equality
\cite[p.~468]{DurrettPTEAuthor2019}.
\index{Hahn partition}\index{signed measure!Hahn partition}
\end{hbdefinition}

\begin{hbtheorem}{thm:p15-w02b-p15-005-hahn-partitions}{proved=annex:p15-w02b-p15-005-hahn-partitions}{Every Hahn partition, and when there is only one}
Let \((A,B)\) be a Hahn partition for a signed measure \(\alpha\).
\begin{enumerate}
\item For a measurable set \(A'\), the pair \((A',A'^c)\) is a Hahn
      partition if and only if \(A\mathbin{\triangle}A'\) is signed-null.
\item The Hahn partition of \(\alpha\) is unique if and only if the
      empty set is the only signed-null measurable set.
\item All Hahn partitions give the same Jordan components: if
      \((A',A'^c)\) is another one, then
      \(\alpha(E\cap A')=\alpha(E\cap A)\) and
      \(\alpha(E\cap A'^c)=\alpha(E\cap B)\) for every measurable \(E\).
\end{enumerate}
\end{hbtheorem}

Part~(1) contains the qualification of \cref{thm:p15-hahn-decomposition}:
the two disagreement pieces there are the two halves of
\(A\mathbin{\triangle}A'\). Part~(2) shows that non-uniqueness is the
typical case. A nonempty signed-null set exists, for instance, whenever
\(\alpha=\int f\,d\mu\) and some nonempty measurable set is
\(\mu\)-null, or whenever \(\alpha\) vanishes identically and \(\Omega\) is
nonempty, since then every measurable set is signed-null.

\begin{hbtheorem}{thm:p15-w02b-p15-005-density-hahn}{proved=annex:p15-w02b-p15-005-density-hahn}{The Hahn positive sets of an integral-defined signed measure}
Let \(\mu\) be a measure on \((\Omega,\mathcal F)\), let \(f\) be an
extended real measurable function with \(\int f^-\,d\mu<\infty\), and
let \(\alpha(E)=\int_Ef\,d\mu\), the signed measure of Example~A.4.1 in
the source and of \cref{thm:p15-signed-measure-examples}.
\begin{enumerate}
\item A measurable set \(N\) is signed-null for \(\alpha\) if and only if
      \(\mu(N\cap\{f\neq0\})=0\).
\item A measurable set \(A\) is a Hahn positive set for \(\alpha\) if
      and only if
      \[
       \{x:f(x)>0\}\subseteq A\subseteq\{x:f(x)\geq0\}
                       \quad\mu\text{-almost everywhere},
      \]
      that is, \(\mu(\{f>0\}\cap A^c)=0\) and
      \(\mu(A\cap\{f<0\})=0\). In particular \(\{f>0\}\) and
      \(\{f\geq0\}\) are Hahn positive sets.
\item Take \(\Omega=[0,1]\) with its Borel sets, \(\mu=\lambda\)
      restricted to them, and \(f(x)=2x-1\). For every Borel
      \(Z\subseteq[0,1]\) with \(\lambda(Z)=0\), both
      \((1/2,1]\cup Z\) and \((1/2,1]\setminus Z\) are Hahn positive
      sets. With \(Z=\{1/2\}\) and with \(Z=C\), the Cantor set of
      \cref{def:p15-cantor-law-construction}, this gives three distinct
      Hahn positive sets \((1/2,1]\), \([1/2,1]\) and
      \((1/2,1]\cup C\). The last contains the point \(0\), where
      \(f(0)=-1<0\).
\end{enumerate}
\end{hbtheorem}

Part~(2) is the source's description
\cite[p.~468]{DurrettPTEAuthor2019}. Part~(3) shows that the
almost-everywhere qualification cannot be dropped: a Hahn positive set
may contain points at which the density is strictly negative, provided
they form a \(\mu\)-null set.

\section{Stieltjes representations and the Jordan components}
\label{sec:p15-w02b-p15-005-stieltjes}

\begin{hbtheorem}{thm:p15-w02b-p15-005-stieltjes-excess}{proved=annex:p15-w02b-p15-005-stieltjes-excess}{The common excess of a Stieltjes representation over the Jordan components}
Let \(\alpha\) be a finite signed Borel measure on \(\mathbb R\), and let
\(F,G\) be bounded, nonnegative, nondecreasing, right-continuous
functions tending to zero at \(-\infty\) with
\[
                     \alpha((-\infty,x])=F(x)-G(x)\qquad(x\in\mathbb R).
\]
Let \(\mu_F,\mu_G\) be their Stieltjes measures
(\cref{thm:p15-stieltjes-construction}), and let \((A,B)\) be a Hahn
partition for \(\alpha\).
\begin{enumerate}
\item \(\alpha(E)=\mu_F(E)-\mu_G(E)\) for every Borel set \(E\).
\item The formula
      \(\rho(E)=\mu_F(E\cap B)+\mu_G(E\cap A)\) defines a finite Borel
      measure, and
      \[
              \mu_F=\alpha^++\rho,\qquad \mu_G=\alpha^-+\rho .
      \]
      It is the only measure with either identity, so it does not
      depend on the Hahn partition. In particular
      \(\mu_F(E)\geq\alpha^+(E)\) and \(\mu_G(E)\geq\alpha^-(E)\)
      for every Borel \(E\), and
      \(F(\infty)+G(\infty)=\lvert\alpha\rvert(\mathbb R)+2\rho(\mathbb R)\),
      where \(F(\infty),G(\infty)\) are the suprema.
\item The following are equivalent: \(\mu_F=\alpha^+\) and
      \(\mu_G=\alpha^-\); \(\rho=0\); \(\mu_F\perp\mu_G\);
      \(F(\infty)+G(\infty)=\lvert\alpha\rvert(\mathbb R)\).
\item Conversely, for every finite Borel measure \(\rho\), the functions
      \[
        x\mapsto\alpha^+((-\infty,x])+\rho((-\infty,x]),\qquad
        x\mapsto\alpha^-((-\infty,x])+\rho((-\infty,x])
      \]
      have all the properties required of \(F,G\) and represent
      \(\alpha\).
\item The choice \(F=G=\mathbf1_{[0,\infty)}\) represents the zero
      signed measure, whose Jordan components are both zero, while
      \(\mu_F=\mu_G=\delta_0\) and \(\rho=\delta_0\). With
      \(F(x)=2\min(\max(x,0),1)\) and \(G(x)=\min(\max(x,0),1)\), the
      represented measure is \(\alpha=\lambda(\,\cdot\,\cap[0,1])\),
      with \(\alpha^+=\alpha\), \(\alpha^-=0\),
      \(\mu_F=2\alpha\), \(\mu_G=\alpha\) and \(\rho=\alpha\).
\end{enumerate}
\end{hbtheorem}

Thus the representations of the cumulative function are exactly the
Jordan pair with one common finite measure added to both terms. The
Jordan pair is the unique representation with least total mass
\(F(\infty)+G(\infty)\), and is singled out among all others by mutual
singularity. This is the precise sense of the qualification in
\cref{thm:p15-signed-cdf-variation} that the two Stieltjes measures
need not be the Jordan components.

\section{Procedure, worked example and limits}
\label{sec:p15-w02b-p15-005-procedure}

\begin{hbmethod}{def:p15-w02b-p15-005-procedure}{Finding every Hahn partition and the Jordan pair}
\begin{enumerate}
\item Produce one Hahn partition \((A,B)\): for an integral-defined
      signed measure take \(A=\{f\geq0\}\)
      (\cref{thm:p15-w02b-p15-005-density-hahn}); in general use
      \cref{thm:p15-hahn-decomposition}.
\item Identify the signed-null sets. For \(\alpha=\int f\,d\mu\) they are
      the measurable sets meeting \(\{f\neq0\}\) in a \(\mu\)-null set.
\item Every other Hahn positive set is \(A\mathbin{\triangle}N\) with
      \(N\) signed-null (\cref{thm:p15-w02b-p15-005-hahn-partitions}).
      The Jordan components \(\alpha^+(E)=\alpha(E\cap A)\) and
      \(\alpha^-(E)=-\alpha(E\cap B)\) are the same for all of them.
\item On the line, given any representing pair \(F,G\), compute
      \(\rho(E)=\mu_F(E\cap B)+\mu_G(E\cap A)\) and subtract its
      cumulative function from both \(F\) and \(G\) to obtain the
      Jordan cumulative functions
      (\cref{thm:p15-w02b-p15-005-stieltjes-excess}).
\end{enumerate}
The side conditions are those of the cited results: the one-sided
convention for \(\alpha\), \(\int f^-\,d\mu<\infty\) for the density
case, and finiteness of \(\alpha\) on the line.
\end{hbmethod}

\begin{hbexample}{def:p15-w02b-p15-005-worked}{Two representations of one signed measure on the unit interval}
Let \(\alpha(E)=\int_E(2x-1)\,d\lambda\) on the Borel subsets of
\([0,1]\), extended to Borel subsets of \(\mathbb R\) by
\(\alpha(E)=\int_{E\cap[0,1]}(2x-1)\,d\lambda\). By
part~(2) of \cref{thm:p15-radon-nikodym-examples}, its Jordan masses
are \(\alpha^+(\mathbb R)=\alpha^-(\mathbb R)=1/4\). Part~(2) of
\cref{thm:p15-w02b-p15-005-density-hahn}, applied on \(\mathbb R\) with
\(\mu=\lambda\) and density \((2x-1)\mathbf1_{[0,1]}(x)\), shows that
\((1/2,1]\), \([1/2,1]\) and \((1/2,1]\cup C\) are all Hahn positive
sets: each contains \((1/2,1]\) and meets \([0,1/2)\) in a
\(\lambda\)-null set. Each gives the same
\(\alpha^+(E)=\int_{E\cap(1/2,1]}(2x-1)\,d\lambda\), because the sets
differ by \(\lambda\)-null sets.

For \(x\in[0,1]\) the cumulative function is
\(\alpha((-\infty,x])=x^2-x\). The pair \(F(x)=x^2\), \(G(x)=x\) on
\([0,1]\), extended by \(0\) to the left and by \(1\) to the right,
represents it: both are bounded, nonnegative, nondecreasing and
continuous. Here \(\mu_F\) has density \(2x\) and \(\mu_G\) has density
\(1\) on \([0,1]\), so with \(A=(1/2,1]\),
\[
 \rho(E)=\int_{E\cap[0,1/2]}2x\,d\lambda+\int_{E\cap(1/2,1]}1\,d\lambda,
 \qquad \rho(\mathbb R)=\frac14+\frac12=\frac34 .
\]
Accordingly \(F(\infty)+G(\infty)=2=\lvert\alpha\rvert(\mathbb R)+2\rho(\mathbb R)
=\tfrac12+\tfrac32\). The measures \(\mu_F,\mu_G\) are not mutually
singular: if \(\mu_F(N)=0\), then \(\lambda(N\cap(0,1])=0\) because the
density \(2x\) is positive there, and so
\(\mu_G(N^c)\geq\lambda((0,1]\setminus N)=1\). By part~(3) of \cref{thm:p15-w02b-p15-005-stieltjes-excess}
this pair is not the Jordan pair. Subtracting
\(\rho((-\infty,x])\) from \(F\) and \(G\) gives the Jordan cumulative
functions \(x\mapsto\alpha^+((-\infty,x])\), which is \(0\) up to
\(1/2\) and \((x-1/2)^2\) on \([1/2,1]\), and
\(x\mapsto\alpha^-((-\infty,x])\), which is \(x-x^2\) on \([0,1/2]\) and
\(1/4\) afterwards.
\end{hbexample}

\begin{hbwarning}{def:p15-w02b-p15-005-limits}{What the non-uniqueness does and does not affect}
The non-uniqueness of a Hahn partition never reaches the Jordan
components or the total variation; it concerns only which signed-null
set is assigned to which side. Statements about the partition itself
are statements up to signed-null sets, and pointwise sign information
about a density holds only \(\mu\)-almost everywhere on a Hahn positive
set. The excess description of
\cref{thm:p15-w02b-p15-005-stieltjes-excess} is proved for finite signed
measures on \(\mathbb R\), the setting of the source's remark; for a
signed measure taking the value \(+\infty\), the cumulative function
\(\alpha((-\infty,x])\) can be infinite and no bounded representing
pair exists.
\end{hbwarning}

The Lebesgue decomposition proof (\cref{thm:p15-lebesgue-decomposition})
uses a Hahn partition of \(\nu_{\mathrm s}-\mu/n\) for each \(n\); by
\cref{thm:p15-w02b-p15-005-hahn-partitions} any choice works, since only
the signed values on subsets of the partition sets are used.
